\scanpage{20}
\mspage{5}
\uline{Proposition} Considérons la cat. des
systèmes $(X,Y,(X_i)_{i\in I}\deleted{\unsure{\}}},\bar T)$ où
$X$ espace top., $Y,(X_i)_{i\in I}$ des sous-espaces fermés,
$\bar T$ voisinage conique de $\unsure{Y|X_i}$
\deleted{compatible}\inserted{compatible} avec
$(X_i)$. (pour les iso), et la catégorie
des systèmes
\[
(X^*,\dot X,(X_i^*),\unsure{Y},
\deleted{\unsure{f}}\dot X\longrightarrow Y)
\]
$X^*,Y$ espaces top., $\dot X,X_i^*$, parties fermées de $X$,
$\dot X\longrightarrow Y$ application continue. Les deux
catégories sont équivalentes.
\inserted{Si $\bar T$ voisinage conique strict}
Pour que $\bar T$
soit \uline{strictement} \deleted{compatible avec les}
\inserted{transversal à} $X_i$,
il faut s. que $\dot X$ soit un sous-espace bord
de $X^*$, \deleted{\unsure{strict}} compatible avec les $X_i^*$.
Pour que $\bar T$ soit équisingulier pour
$(X_i)$, il faut s. que
\diagram{\[
\begin{tikzcd}
\dot X\arrow[r]&Y\\
(\dot X_i)\arrow[u,hook]&
\end{tikzcd}
\]}
soit une fibration localement triviale
simultanée (où on pose $\dot X_i=X_i\cap\dot X$)
\margin{19}
